% sections/overview.tex -- Figure 1 (F1): schematic of the method.
% \input at the end of sections/intro.tex; the TikZ source is figures/f1_method.tex.
\begin{figure}[t]
  \centering
  % figures/f1_method.tex -- Figure 1: the five parts of the method (TikZ,
% grey-scale only). \input inside the figure environment of
% sections/overview.tex; uses the notation of macros.tex. Solid arrows: the
% update path, run for every atomic edit. Dashed arrows: the rebuild path,
% which affects only the cost. The algebra of a deletion is in Figure 2
% (figures/f2_deletion.tex), the Morse order in Figure 3 (figures/f3_morse.tex).
\begin{tikzpicture}[
    font=\footnotesize,
    >={Stealth[length=2mm,width=1.6mm]},
    part/.style={draw=black!80, rounded corners=2pt, fill=white, align=flush left,
      inner sep=4pt, text width=#1},
    core/.style={part=#1, fill=black!8},
    num/.style={circle, fill=black, text=white, inner sep=0pt, minimum size=3.8mm,
      font=\scriptsize\bfseries},
    flow/.style={->, semithick},
    rb/.style={->, semithick, dash pattern=on 3pt off 2pt},
    lbl/.style={font=\scriptsize, inner sep=1.5pt},
  ]
  % ----- update path ----------------------------------------------------------
  \node[part=3.25cm] (edits) at (0,0) {%
    \hspace*{2.4mm}\textbf{Atomic edits}\\
    an update $u_t$ is a batch of insertions $\Ins(\sigma)$ and deletions
    $\Del(\sigma)$, applied one at a time};
  \node[core=5.3cm, right=8mm of edits] (drb) {%
    \hspace*{2.4mm}\textbf{Dynamic reduced basis}\\
    per dimension: slots $(R_j,V_j)$, $R=\Bm{k}V$ reduced, $V$ invertible\\
    $\Ins$: reduce the new column against the pivots\\
    $\Del$: rewrite the footprint $\occ(\sigma)$, Case~A or~B
    (\cref{fig:deletion})};
  \node[part=3.75cm, right=8mm of drb] (betti) {%
    \hspace*{2.4mm}\textbf{Exact Betti numbers}\\
    $\beta_k=n_k-r_k-r_{k+1}$ with $r_k=\#\{j:R_j\ne0\}$, after every edit and
    for every order of the cells};
  % ----- rebuild path -----------------------------------------------------------
  \node[part=5.3cm, below=10mm of drb] (rebuild) {%
    \hspace*{2.4mm}\textbf{Optional rebuild}\\
    at initialisation, and after $u_t$ if $m$ updates have passed or
    $\nnz(R)+\nnz(V)$ has grown by a factor $1+\alpha$: static reduction with
    clearing, tracking $V$};
  \node[part=3.25cm, anchor=north east] (morse) at (edits.east |- rebuild.north) {%
    \hspace*{2.4mm}\textbf{Morse-induced order}\\
    keys from a collapse matching; additions only in critical columns
    (\cref{fig:morse})};
  % ----- numbers --------------------------------------------------------------
  \foreach \b/\n in {edits/1, drb/2, betti/3, rebuild/4, morse/5}
    \node[num] at (\b.north west) {\n};
  % ----- arrows ---------------------------------------------------------------
  \draw[flow] (edits) -- (drb);
  \draw[flow] (drb) -- (betti);
  \draw[rb] (betti.south) |- node[lbl, pos=0.25, right] {after $u_t$} (rebuild.east);
  \draw[rb] (morse.east |- rebuild.west) -- node[lbl, above] {keys} (rebuild.west);
  \draw[rb] (rebuild.north) -- node[lbl, right] {fresh basis} (drb.south);
\end{tikzpicture}

  \caption{The method. \emph{Top:} an update is processed one atomic edit at a
    time on the dynamic reduced basis of each boundary matrix
    (\cref{sec:drb}). An insertion reduces the new column against the pivots; a
    deletion rewrites every slot in the footprint $\occ(\sigma)$ of the deleted
    simplex, by a cycle (Case~A) or by the slot with the smallest low (Case~B).
    The Betti numbers are exact after every edit, whatever the order of the
    cells. \emph{Bottom, dashed:} a rebuild replaces the order by the Morse
    order of a collapse matching and recomputes the basis by the static
    reduction with clearing (\cref{sec:morse}). Under that order every gradient
    pair is an apparent pair and all column additions happen in columns of
    critical cells. Between rebuilds the Morse order is only a heuristic for
    sparsity (\cref{rem:morse-scope}). The ablations of \cref{sec:exp-methods}
    change the bottom path: \method{drb-id-rb} rebuilds in insertion order, and
    \method{drb-id} and \method{mmhm-norb} never rebuild.}
  \label{fig:method}
\end{figure}
